\section{总结与延伸：Love the Transformer，但不要停止追问}

这堂课的标题不是无条件崇拜 Transformer，而是接受一种研究事实：它把 content-based routing、parallel representation learning、residual computation 和 scalable matrix operations 结合成极强 substrate，并推动 NLP 从专用 pipeline 走向共享系统。真正值得“love”的不是某个固定 block，而是 consolidation、learned interface 与 algorithm--system co-design 的方法。

\subsection{九条可迁移的工程结论}

下面的结论可用于评估新的 sequence model、agent architecture 或 long-context 系统。它们把历史、机制和系统证据连接起来，避免团队只比较参数量或单个 benchmark。

\begin{enumerate}
  \item \textbf{先问系统吸收了哪些接口。} Consolidation 的价值来自减少 hand-written boundaries，而不是模型名称更新。
  \item \textbf{区分 expressivity 与 trainability。} LSTM 能表达复杂算法，仍可能因 sequential critical path 不适合 scale。
  \item \textbf{把 attention 视为 content routing。} Q/K 决定连接，V 决定内容，FFN 与 residual 决定后续变换。
  \item \textbf{记录失败的原始目标。} Parallel reading 成功，parallel writing 仍受 mode/order 限制。
  \item \textbf{Position 是 pairwise interface。} 长度外推取决于 encoding、frequency、training range 与 attention behavior。
  \item \textbf{复杂度必须下钻到 data movement。} FLOPs、activation、bandwidth、kernel regularity 和 latency 要分开报告。
  \item \textbf{模型内外能力要合理分工。} Exact、fresh、可验证操作适合工具，语义组合与策略选择适合模型。
  \item \textbf{多 agent 必须有 decomposition、coordination、verification。} 角色提示不能代替系统协议。
  \item \textbf{用范围声明保护结论。} 一张 attention map、一段音乐 demo 或一个 benchmark 都只支持有限主张。
\end{enumerate}

\begin{importantbox}{把本讲压缩成一句话}
Transformer 的持久价值不在于“attention 永远是最终答案”，而在于它展示了如何把\textbf{通用问题表达、内容寻址、并行矩阵计算、可训练残差结构与硬件友好实现}组合成可扩展 substrate；下一代架构仍必须回答这些约束。
\end{importantbox}

\subsection{一张系统审计表}

\begin{longtable}{@{}L{0.18\textwidth}L{0.35\textwidth}L{0.35\textwidth}@{}}
\toprule
维度 & 应追问的问题 & 常见误读 \\
\midrule
Representation & token、position、state 如何表示？ & 统一 tokenization 等于统一语义 \\
Routing & 信息按内容、位置还是固定 topology 交换？ & attention weight 就是完整解释 \\
Dependency & 哪些步骤必须 sequential？ & 训练可并行等于 inference 可并行 \\
Complexity & FLOPs、activation、bandwidth 各是多少？ & Big-O 更低就一定更快 \\
Memory & KV、外部 retrieval、长期状态在哪里？ & context 能装下就能有效使用 \\
Interface & 输出由模型生成还是调用工具？ & 工具越多能力越可靠 \\
Learning & 梯度、in-context memory、反馈分别更新什么？ & 所有“学习”都等同于参数训练 \\
Evidence & 结果、demo、interpretability 各证明什么？ & 局部证据被外推为通用能力 \\
\bottomrule
\end{longtable}

\subsection{拓展阅读}

\begingroup
\footnotesize
\begin{itemize}[itemsep=0pt,topsep=0.15em,parsep=0pt]
  \item Vaswani et al., \href{https://arxiv.org/abs/1706.03762}{\emph{Attention Is All You Need}}
  \item Bahdanau, Cho, and Bengio, \href{https://arxiv.org/abs/1409.0473}{\emph{Neural Machine Translation by Jointly Learning to Align and Translate}}
  \item Gehring et al., \href{https://arxiv.org/abs/1705.03122}{\emph{Convolutional Sequence to Sequence Learning}}
  \item Shaw, Uszkoreit, and Vaswani, \href{https://arxiv.org/abs/1803.02155}{\emph{Self-Attention with Relative Position Representations}}
  \item Huang et al., \href{https://arxiv.org/abs/1809.04281}{\emph{Music Transformer}}
  \item Roy et al., \href{https://arxiv.org/abs/2003.05997}{\emph{Efficient Content-Based Sparse Attention with Routing Transformers}}
  \item Ainslie et al., \href{https://arxiv.org/abs/2305.13245}{\emph{GQA: Training Generalized Multi-Query Transformer Models}}
  \item Dao et al., \href{https://arxiv.org/abs/2205.14135}{\emph{FlashAttention}}
  \item Schick et al., \href{https://arxiv.org/abs/2302.04761}{\emph{Toolformer}}
\end{itemize}

阅读时建议沿用系统审计表，并持续区分三层：论文证明的 mechanism/result、课堂给出的 historical/engineering judgment、以及今天读者可能做出的 extension。只要一条结论跨越了 source boundary，就应明确标注它是推论而非课堂原话。
\endgroup

\end{document}
